%\documentclass[aps,rmp,groupedaddress,twocolumn,10pt]{revtex4-2}
\documentclass[aps,prd,reprint,superscriptaddress]{revtex4-2}

\usepackage{graphicx}
\usepackage{hyperref}
\usepackage{amssymb}
\usepackage{amsmath}
\usepackage[german,english]{babel}

\begin{document}

\title{Cosmological Relativity: Resolving the Isotropy and Cosmic Time Problems \\in Precision Cosmology}

\author{Daryl Janzen}
\email[]{daryl.janzen@usask.ca}
\affiliation{Department of Physics and Engineering Physics, University of Saskatchewan, 116 Science Place, Saskatoon, SK S7N 5E2, Canada}

\date{\today}

\begin{abstract}
The precise isotropy of the CMB monopole's redshift, despite the highly anisotropic matter distribution in our observable universe, presents an unresolved crisis in precision cosmology. A conservative estimate based on standard perturbation theory and the central limit theorem predicts redshift anisotropies at the $\delta z/z\gtrsim1\%$ level---large enough to overwhelm the primordial CMB anisotropies and even the dipole. That no such effect is observed is not merely a fine-tuning issue; it is an empirical failure of GR-based reasoning, which overlooked an effect that, had it been considered, would have led to a falsified prediction. The absence of this expected redshift contrast requires an unexplained suppression of more than three orders of magnitude, demanding a theoretical resolution. Cosmological Relativity (CR) provides a natural resolution to this isotropy problem, as well as to several long-standing epistemic failures of operational relativity revealed by modern cosmology. CR is a geometric augmentation of GR that retains its exact empirical success in local space-time while resolving fundamental cosmological tensions. Unlike a mere conceptual reinterpretation, CR introduces new mathematical structure---a preferred foliation with absolute cosmic time and diffeomorphic flexibility between real space and cosmological space-time slices---allowing it to accommodate observed isotropy while preserving local covariance. This distinct geometric framework leads to new predictions for emergent features of our universe, including the nature of black holes and the phenomenological flat $\Lambda$CDM expansion rate derived from first principles. By resolving the core epistemic problems of FLRW (GR+RW) cosmology---along with paradoxes related to time travel and black holes---CR offers a more complete foundation for modern physics while remaining fully consistent with established observational constraints. It also provides a clear empirical prediction with the potential to resolve the Hubble tension. CR follows the historically established progression of scientific theory, in which phenomenological descriptions evolve to distinguish observed phenomena from underlying physical reality once fundamental inconsistencies in naive and superficial interpretations have been identified.
\end{abstract}

%\keywords{Relativity theory,Cosmic time,Block universe,Black hole physics,Cosmic microwave background (CMB),Gravitational collapse,Ontology of existence,Noumenology and phenomenology,Philosophy of science,General relativity}

\maketitle

\section{Introduction}\label{sec:intro}

\subsection{General relativity and its limitations in cosmology}

General relativity (GR) revolutionized our understanding of gravity, describing it not as a force but as the curvature of space-time caused by mass-energy. In this framework, space-time is a four-dimensional continuum where the metric $g_{\mu\nu}$ governs the relationships between events, determining how objects move and how light propagates. The Einstein field equations connect this curvature to the stress-energy tensor, providing a local description of how matter and energy shape the geometry of space-time  \cite{Einstein1915,Einstein1952}. GR's predictions, from gravitational waves to black holes, have been confirmed with remarkable precision. However, the theory faces conceptual and technical challenges, particularly regarding singularities and its global application to cosmology.

To address these challenges, the Arnowitt-Deser-Misner (ADM) formalism reformulates GR by splitting space-time into a series of spatial hypersurfaces that evolve over time \cite{Arnowitt1959, Arnowitt1962}. This decomposition expresses the dynamics of space-time in terms of the spatial metric $h_{ij}$, the lapse function $N(x)$, and the shift vector $N^i(x)$, allowing the Einstein field equations to be recast as evolution equations for these quantities. ADM provides a powerful computational framework for numerical relativity, facilitating the analysis of gravitational wave propagation, black hole mergers, and other dynamic systems \cite{Baumgarte2010}. While it encodes GR's space-time geometry mathematically, it does not introduce additional physical structure. 

While ADM excels in practical applications, its reliance on foliated hypersurfaces reveals deeper limitations in GR itself. Near singularities, where curvature diverges, the mathematical coherence of both GR and ADM breaks down, reflecting the difficulty of handling such regions in a physically meaningful way. This limitation underscores a broader issue: GR's reliance on a rigid, globally defined space-time metric constrains its ability to account for certain evolving structures of the observable universe. Nevertheless, ADM remains indispensable in numerical studies, offering a highly convenient tool for exploring space-time dynamics under complex conditions.

In addition to these challenges, cosmology has revealed significant tensions with GR-based expectations. In particular, there are six epistemically equivalent problems that have received mixed appreciation: the expansion/existence problem, the flatness problem, the horizon problem, the dark matter/dark energy/cosmological constant problem, the isotropy problem, and the cosmic time problem. The epistemic equivalence of these problems lies in the fact that while the standard cosmological model successfully \textit{describes} a universe exhibiting these properties, each of them highlights a fundamental tension between empirical constraints and core principles of GR. Specifically, GR not only fails to provide a theoretical basis for why the universe \textit{should} exhibit these properties, but more problematically, it leads to the expectation that the universe \textit{should not} exhibit them. We'll start with the first four of these problems, which have been more widely recognized in the literature.

\subsection{The classic cosmological problems and their theoretical tensions}

The expansion/existence problem was most famously appreciated, separately, by Eddington \cite{Eddington1933} and Hoyle. 
To understand the problem, we start by noting that a universe containing matter (with present energy density $\rho_{m,0}$), radiation (with present energy density $\rho_{r,0}$), and a cosmological constant $\Lambda$---corresponding to an equivalent constant vacuum energy density, or ``dark energy,'' given by $\rho_{\Lambda} = \frac{\Lambda}{8\pi G}$---will accelerate as
\begin{equation}
    \ddot{a}=-\frac{4\pi G}{3}\left(2\rho_{r,0}\,a^{-3} + \rho_{m,0}\,a^{-2}\right) + \frac{\Lambda}{3}\,a,\label{eq:accel}
\end{equation}
where $a=a(t)$ is the scale-factor which describes cosmic expansion according to the Robertson-Walker (RW) line-element,
\begin{equation}
    ds^2=-dt^2+a(t)^2\left(\frac{1}{1-kr^2}dr^2+r^2d\Omega^2\right),\label{eq:RW}
\end{equation}
where the allowed values $k=-1,\,0,\,1$ describe spaces that are, respectively, hyperbolic, flat, or spherical. The expansion rate is typically characterized by the Hubble parameter, $H = \dot{a}/a$, which describes the relative rate of expansion at any given time. The present-day value, $H_0$, known as the Hubble constant, sets the observed expansion rate of the universe today. Another key quantity is the total density parameter, $\Omega = \rho / \rho_c$, where $\rho_c = 3H^2 / 8\pi G$ is the critical density required for a spatially flat universe, and is determined by solving Einstein's field equations with the RW metric. 

General relativistic universes that are well-approximated by Eq.~(\ref{eq:RW}), as ours is understood to be when considered on large enough scales, which contain only matter, radiation, and a cosmological constant, accelerate according to Eq.~(\ref{eq:accel}). We refer to these GR+RW model universes as Friedmann-Lema\^{i}tre-Robertson-Walker (FLRW) models, after their primary contributors \cite{Friedmann1922,Friedmann1924,Lemaitre1927,Robertson1935,Walker1937}. A \textit{big bang} FLRW universe---which expands from an initial singularity at $a(0)=0$---begins, according to Eq.~(\ref{eq:accel}), with matter and radiation driving deceleration at an initially infinite rate. Thus, both the apparent expansion of space \cite{Hubble1929} and the very existence of the universe are treated as mere consequences of initial conditions at the singularity, rather than natural predictions of the model. The singularity at $a=0$ also corresponds to the breakdown of classical GR \cite{Hawking1970}, reinforcing that the model offers no explanation for the initial state of the universe. Therefore, neither the property that our universe expands, nor the property that it even exists, is \textit{explained}. In fact, the model indicates that all the matter and radiation in the universe initially works \textit{against} this expansion while the contribution of the cosmological constant is initially negligible; that it is only through significant inertia imparted at an initial moment when the physics breaks down, that such expanding universes exist.

While inflation is often invoked to resolve various cosmological problems, it does not address the fundamental issue here: FLRW models, even in the presence of inflationary fields, still rely on an initial state where classical GR breaks down, and expansion is an unexplained consequence of initial conditions rather than a natural outcome of the model. Additionally, appealing to quantum gravity does not mitigate the issue, as no complete theory of quantum gravity exists, and any such theory would still need to provide a transition to classical FLRW behavior that accounts for the apparent inertial legacy of expansion without assuming it as an arbitrary initial condition. The expansion/existence problem is thus not merely a question of unknown early-universe physics, but rather a deeper issue regarding the explanatory power of the standard framework.

Both the flatness problem and the horizon problem are widely appreciated, as is their favoured resolution, the inflation scenario \cite{Guth1981,Linde1982,Linde1983,Albrecht1982}. The flatness problem arises because FLRW universes permit a wide range of possible spatial curvatures, yet observations indicate that the universe is nearly exactly flat. However, a precisely flat universe is a measure-zero solution within the space of possible initial conditions---any small deviation from flatness in the early universe should have been exponentially amplified over cosmic time so that the universe should now have significant curvature. Similarly, the horizon problem stems from the extreme uniformity of the cosmic microwave background (CMB) temperature across regions that, under standard GR evolution, should never have been causally connected. Any two points separated by more than $\sim 2$ degrees on the CMB sky should have evolved independently, yet the observed temperature fluctuations are correlated to better than 1 part in $10^5$, despite there being no causal mechanism to enforce this. Inflation is widely considered the best available resolution, as an early phase of exponential expansion would dynamically drive $\Omega$ toward unity while stretching once---causally connected regions far apart. However, inflation introduces its own theoretical challenges---including the measure problem, fine-tuning of the inflationary potential, and the apparent inevitability of eternal inflation---raising concerns over whether it constitutes a genuine predictive framework or merely a parameter-tuning mechanism \cite{Ijjas2014, Ijjas2017}. Nevertheless, both problems are epistemically equivalent to the expansion/existence problem: the standard model describes a universe bearing the empirically constrained properties without issue, but the observed properties themselves lack theoretical motivation without an additional mechanism to drive the universe's evolution so that it would match what we observe.

The energy budget of the universe presents an even deeper problem. While standard cosmology describes the evolution of the universe using precisely measured parameters, it does not explain why these parameters take the values they do---or even what most of the universe is made of. Observations indicate that baryonic matter and radiation account for only $\sim4\%$ of the total energy density of the universe, while the remaining $\sim96\%$ is attributed to invisible dark matter and dark energy components \cite{Zwicky1933,Rubin1980,Riess1998,Perlmutter1999,Aghanim2020}. Furthermore, theoretical estimates of vacuum energy from quantum field theory exceed the observed dark energy density by an astounding 120 orders of magnitude \cite{Weinberg1989,Carroll1992}, yet for reasons unknown, this vacuum energy appears to have no observable effect on expansion. While empirical parameter fitting provides a consistent description, the theoretical justification for these components remains elusive. The difficulty is compounded by the so-called Hubble tension, wherein independent measurements of the Hubble constant $H_0$ disagree at a statistically significant level ($>5\sigma$) \cite{Freedman2017,Verde2019,Riess2024}, raising concerns that the inferred energy budget may not even be internally consistent. While proposals such as early dark energy \cite{Kamionkowski2023} seek to resolve this discrepancy by varying the vacuum energy from that of a cosmological constant, they introduce additional degrees of freedom and fine-tuning requirements, reinforcing the broader issue: while the standard model \textit{describes} the observed expansion history well, it does not \textit{explain} why its parameters take the values they do---nor even \textit{what} dark matter and dark energy are. Thus, the ``energy budget problem''---like the flatness and horizon problems---is another example of the standard model accurately fitting the data, yet failing to justify its own assumptions.

\subsection{The isotropy and cosmic time problems}

Neither the isotropy problem nor the cosmic time problem has received the attention they are due, as these problems are epistemically equivalent to, and as significant as, the four problems noted above. The isotropy problem is distinct from, yet closely related to, the horizon problem. While the horizon problem concerns the universe's isotropy \textit{at the time of recombination}, the isotropy problem is concerned with the evidence that cosmic expansion \textit{since} recombination has remained perfectly isotropic despite significant anisotropies in the matter distribution. The isotropy problem refers to a tension between the anisotropy of the observable universe's matter distribution and the isotropy of the CMB monopole's overall \textit{redshift}, which together indicate that our significantly lumpy, inhomogeneous universe has expanded uniformly at a continuously variable expansion rate \textit{throughout} the 13.8 billion years since recombination. The problem is not that the high degree of isotropy of the CMB monopole indicates the universe was nearly homogeneous \textit{at the time} of recombination in spite of prior causal disconnection; it's that in order for that isotropy to have been \textit{preserved} and still be apparent after 13.8 billion years of travel through expanding space, the cumulative expansion over that period must have been \textit{precisely isotropic}---meaning that in every direction we look, at every point along each photon's path, \textit{despite} significant anisotropy in the distribution of matter---which is \textit{supposed to} affect expansion locally according to GR---the cosmic expansion rate seems to have been precisely the same.

CMB photons that have travelled unimpeded through the universe since recombination are not observed in the same state they were emitted; their redshifts arise from two distinct effects: the expansion of space and inhomogeneous gravitational fields. Cosmological redshift is not an intrinsic property of the emitting source but instead results from the continuous stretching of photon wavelengths as space itself expands along their entire trajectory. In contrast to Doppler and gravitational redshifts, cosmological redshift is not about here and now relative to there and then, but about the continuously varying path of the photon itself. This means that the cosmological redshift of a photon is not fixed at the moment of emission, but \textit{accumulates progressively} as the photon propagates through the expanding universe. 

The Sachs-Wolfe (SW) effect \cite{Sachs1967} results from gravitational redshift at recombination, where photons originating from overdense regions experience a slight redshift, while those from underdense regions experience a slight blueshift, imprinting anisotropies with peak fluctuations at $\Delta T/T \approx 3 \times 10^{-5}$ on large angular scales and finer-scale variations at $\Delta T/T \approx 3 \times 10^{-6}$ \cite{Smoot1992, Bennett2003, Aghanim2020b}. Additional secondary effects---such as the integrated SW effect, gravitational lensing, the Rees-Sciama effect, and the Sunyaev-Zel'dovich effect---modify the CMB signature along its 13.8-Gyr journey through expanding space \cite{Ade2016, Lewis2006, Rees1968, Crittenden1996, Sunyaev1970, Sunyaev1972}. \textit{Unlike cosmological redshift, which is continuously acquired as a direct consequence of space itself expanding along the photon's trajectory, gravitational redshift effects result from localized variations in potential and are acquired at specific regions along the photon's path.}

Since cosmological redshift depends on the detailed expansion history along the entire photon trajectory, \textit{any anisotropies in the expansion rate should have induced corresponding anisotropies in redshift.} The observed $z=1089$ expansion redshift represents a large-scale average, but it does not preclude local deviations. Given the significant anisotropy in the matter distribution, GR predicts that expansion rates should have varied slightly along different lines of sight, introducing measurable $\delta z/z$ perturbations, which would manifest as corresponding $\delta T/T$ anisotropies in the CMB. In fact, general-relativistic expansion anisotropies have been explored in the literature in attempts to explain the apparent acceleration of cosmic expansion \cite{Wiltshire2009, Seifert2025}. Since structure formation has led to significant inhomogeneity in the matter distribution over cosmic time, \textit{GR-based intuition suggests that observed density variations should have influenced the total amount of expansion experienced by CMB photons.} If so, the expected anisotropies should be measurable, and their magnitude can be estimated based on the statistical properties of matter inhomogeneity along a typical CMB photon's path.

In order to estimate the expected $\delta T/T=\delta z/z$ contrast due to expansion anisotropy arising from matter density inhomogeneity, we choose 8-$h^{-1}$Mpc as a scale at which matter distributions locally tend to become roughly uniform, which still captures the large-scale statistical distribution in matter inhomogeneity through the cosmological parameter $\sigma_8(z)$, where $\sigma_8(0)\approx0.8$. Since Friedmann's equations describe the instantaneous expansion rate of any region with approximately uniform density, they should describe the expansion rate along a CMB photon's path with slowly varying density. In particular, these equations tell us 
\begin{equation}
    \frac{d\rho}{\rho}\sim -\frac{da}{a}.
\end{equation}
Thus, density perturbations are expected to lead to similar magnitude scale-factor perturbations, with an increase in overall density affecting a decrease in the expansion of space. Furthermore, $1+z=a_{\mathrm{ob}}/a_{\mathrm{em}}$ means that two photons emitted at $a_{\mathrm{em}}=a_{\mathrm{lss}}$ and observed at $a_{\mathrm{obs}}$ and $a_{\mathrm{obs}}+da$, we have
\begin{equation}
    \frac{da}{a}\sim\frac{dz}{z}.
\end{equation}
Essentially, what this means is that when we look back to the last scattering surface, in a direction where the matter density has been overall slightly lower along a CMB photon's path, space should have expanded slightly more and the overall redshift should therefore be slightly greater. Since the overall distance the photon travelled still must be 13.8~Gly, such photons must have originated from a slightly closer comoving distance than the average. Similarly, in a slightly overdense direction, space will not have expanded as much overall since recombination, so the expansion redshift will be less while photons will have travelled further. 

In the extreme limit, the universe might have been so overdense as to have halted expansion in one direction, in which photons from the last scattering surface will have travelled 13.8~Gly through static space resulting in zero expansion redshift, while in another direction the universe may have been expanding so rapidly as to induce an extraordinarily large redshift in photons from the last scattering surface that may have been emitted from relatively nearby comoving coordinates. With more modest anisotropies, within an inhomogeneously expanding universe the last scattering surface would be an approximately spherical, but nevertheless distorted 2-dimensional surface surrounding our location, with distances and associated apparent CMB redshifts varying based on the density perturbations and associated expansion perturbations encountered along the way.

To estimate the degree of redshift contrast we should observe due to density perturbations and their expected influence on local expansion rates, there are several things we need to account for at once. Start by considering the path of CMB photons in a particular direction as a tube whose comoving diameter is currently $8~h^{-1}$Mpc, and whose length is the present proper distance to the last scattering surface, $d_{\mathrm{lss},0}$. Assuming small perturbations whose effects to the local expansion rate roughly cancel, the standard model gives the approximate value, $d_{\mathrm{lss},0}\approx9.42~h^{-1}\mathrm{Gpc}$ \cite{Ryden2017}. Since $d_{\mathrm{lss},0}\gg8~h^{-1}$Mpc, rather than treating density perturbations as continuously varying across the length of this tube, we can apply the central limit theorem to estimate the standard deviation across discrete blocks of length $8~h^{-1}$Mpc whose density contrasts are given by $\sigma_8(z)$. In order to do this, we first need to calculate
\begin{equation}
    \sigma_{8,\mathrm{eff}}=\frac{\int_0^{d_{\mathrm{lss},0}}\sigma_{8}(z)dr}{\int_0^{d_{\mathrm{lss},0}}dr}=\frac{\int_0^{z_{\mathrm{lss}}}\sigma_{8}(z)H(z)^{-1}dz}{\int_0^{z_{\mathrm{lss}}}H(z)^{-1}dz},\label{eq:sigma_eff}
\end{equation}
the effective standard deviation along this tube, which accounts both for the diminishing value of $\sigma_8(z)$ as density variations effectively drop to $\sigma_8(z_{\mathrm{lss}})\approx0$ at $z_{\mathrm{lss}}=1089$, and the fact that cosmic expansion has meant CMB photons crossed high-$z$ regions more rapidly, so the effect of low-$z$ density perturbations is relatively amplified. For our purposes, we use the approximate flat $\Lambda$CDM expansion rate 
\begin{equation}
    H(z)=H_0\sqrt{\Omega_m(1+z)^3+\Omega_{\Lambda}} \label{eq:H}
\end{equation}
to capture the large-scale average. Finally, to capture the variation in $\sigma_8(z)$, we write $\sigma_8(z)=\sigma_8(0)D(z)$, where
\begin{equation}
    D(z)=\frac{H(z)}{H_0}\int_z^{\infty}\frac{(1+z')dz'}{H(z')^3}\left[\int_0^{\infty}\frac{(1+z')dz'}{H(z')^3}\right]^{-1}\label{eq:Dz}
\end{equation}
is the linear growth factor (normalised to $D(0)=1$) for a flat $\Lambda$CDM universe from standard perturbation theory \cite{Peebles1980}. 

Intriguingly, for the Planck best-fit value of \( \Omega_m = 0.315 \), the normalisation factor in Eq.~(\ref{eq:Dz}),  
\begin{equation}
    I = \int_0^{\infty} \frac{(1+z')dz'}{H(z')^3},
\end{equation}
evaluates to \textit{precisely} 1 when using the flat $\Lambda$CDM expansion rate. However, $I$ is an exponential decay function that varies significantly with $\Omega_m$. There is no theoretical reason why this integral should take a value of precisely 1 at $\Omega_m=0.315$, yet this is coincidentally the value empirically inferred from analyses that depend on this very normalisation factor. If this factor were unintentionally omitted or implicitly assumed to be 1 in analysis pipelines, it would artificially drive analysis pipelines towards $\Omega_m=0.315$. This coincidence may warrant further investigation, as it raises the possibility of a subtle but systematic bias in standard parameter estimation.

Inserting Eqs.~(\ref{eq:H})-(\ref{eq:Dz}) into Eq.~(\ref{eq:sigma_eff}) and evaluating for $\sigma_8(0)=0.8$, $z_{\mathrm{lss}}=1089$, $\Omega_m=0.315$ and $\Omega_{\Lambda}=1-\Omega_m$ yields $\sigma_{8,\mathrm{eff}}=0.285$.\footnote{As a check that this value of $\sigma_{8,\mathrm{eff}}$ is not dominated by low-$z$ fluctuations where $8~h^{-1}$Mpc density perturbations cover wide angles and there are not a statistically significant number of lines of sight, we note that since $\sigma_8(z=1)\approx0.5$ remains large, excluding the local universe and integrating over $z=[1,1089]$ still yields $\sigma_{8,\mathrm{eff}}=0.17$, confirming that our result is not dominated by low-$z$ fluctuations.} Now, the calculation of the density anisotropy along this path, and the associated redshift contrast across separate lines of sight, by applying the central limit theorem. While large-scale structures are significantly correlated above the $8~h^{-1}$Mpc-scale, it is sufficient to \textit{under}estimate the density contrast by assuming no correlation along the comoving tube that models the photon's path, noting that any larger structure correlations (e.g.\ within basins of attraction of mean length $138~h^{-1}$Mpc \cite{Valade2024}) would only decrease the number of effective bins in the central limit. Thus, looking outwards we draw randomly from spherical shells of thickness $8~h^{-1}$Mpc, whose density contrast is given by $\delta\rho/\rho=\sigma_{8,\mathrm{eff}}$. Thus, the total density contrast along this path, and the corresponding temperature contrast due to expansion anisotropies, are given by
\begin{equation}
    \frac{\delta T}{T}=\frac{\delta z}{z}\approx\frac{\delta\rho}{\rho}\equiv\sigma_{\mathrm{path}}\approx\frac{\sigma_{8,\mathrm{eff}}}{\sqrt{N}}\approx 8\times10^{-3},\label{eq:dz}
\end{equation}
where $N=9.42~h^{-1}\mathrm{Gpc}/8~h^{-1}\mathrm{Mpc}\approx1200$ is the number of independent and identically distributed $8~h^{-1}$Mpc-bins along a radial line of sight to the surface of last scattering.

In light of the unexpectedly large redshift variation predicted in our highly conservative estimate, Eq.~(\ref{eq:dz}), it is crucial to reconsider the extent to which the FLRW model's fundamental assumptions are justified. The RW line-element is first formulated on kinematical grounds, defining the large-scale structure of space-time, before being made dynamical by solving the Einstein field equations, at which point the energy density and pressure are identified as those of a perfect fluid. Thus, the assumptions defining the line-element itself---in conjunction with the assumption that GR is the valid governing theory---constitute the core hypotheses of modern cosmology.

The RW line-element is constructed on four key assumptions, all of which were historically justified only as approximate large-scale averages rather than exact principles \cite{Robertson1933}. The first assumption is that the mean motion of galaxies defines a coherent bundle of non-intersecting geodesics with a common origin, along which equal proper times define a preferred simultaneity. Second, this preferred cosmic time is assumed to be always orthogonal to the spatial hypersurfaces defined by its constant values, so that `simultaneous' events are also synchronous in the cosmic rest frame. Having introduced this special foliation of space-time, we then assume that at sufficiently large scales, space is isotropic and homogeneous, and that the universe undergoes a uniform, isotropic expansion in cosmic time, consistent with the large-scale applicability of Hubble's law. These assumptions lead directly to the RW line-element, Eq.~(\ref{eq:RW}).

In a sense, our estimated redshift contrast, which is based on the observed density contrast in our universe and the size of the observational universe inferred from standard cosmological parameters, verifies the core assumptions of the FLRW model to the extent that they were originally supposed to hold. While this is a conservative estimate, the prediction that CMB redshifts should vary by $\sim 1\%$ across the sky---corresponding to a temperature anisotropy at the $10^{-2}$ level---demonstrates that the model's original assumptions were at best valid as large-scale approximations. However, the fact that the observed power spectrum constrains $\delta T/T$ to orders of magnitude smaller values presents a fundamental contradiction.

The problem, however, is that the standard CMB analysis never accounted for this effect. The theory assumes in principle that the effect of cosmic expansion on CMB redshifts is \textit{perfectly} isotropic; it ignores the effect of density perturbations to the cosmic expansion rate completely and assumes a perfectly isotropic, homogeneous, and uniformly expanding universe, as described by the RW line-element and the Friedmann equations. Density parameters are modelled on an assumption that the angular diameter distance to the last scattering surface is perfectly isotropic---and indeed, we see no distortion effects in the CMB that would be explained by a surface that is actually deformed due to variable expansion rates encountered along the way. The standard model accounts only for secondary effects, such as gravitational redshift and lensing, that arise from inhomogeneous structures \textit{in} space. However, it neglects the primary effect of inhomogeneous expansion \textit{of} space itself. This omission is critical: had the missing expansion anisotropies been accounted for, the resulting theoretical prediction would have been starkly incompatible with observations. 

The problem is that redshift anisotropies that should have been an unavoidable consequence of observed density anisotropies in our universe were never accounted for. Had they been included in standard treatments, they would have led to a falsified prediction decades ago. Instead, the standard model accidentally ignored an effect that, had it been considered, would have directly contradicted observations. These anisotropies, which are beautifully explained by structure formation within a precisely isotropic expansion framework, would be overwhelmed by the redshift contrast of order $\sim1\%$ expected from inhomogeneity-driven expansion anisotropies. The fact that no such effect is seen in the data suggests that the underlying assumption of a statistically isotropic expansion is not merely an approximation, but an exact empirical constraint---one that demands an explanation beyond standard GR-based expectations.

The standard model has not merely assumed uniform expansion---it has accidentally validated a theoretical assumption that should have been questioned from the outset. This extreme precision is not a triumph of prediction, but a consequence of ignoring an effect that, had it been considered, would have led to a direct contradiction with observation. The expansion history of our universe, looking in \textit{any} two directions across the sky, has been \textit{precisely} the same, going all the way back to recombination. The fact that the expected expansion anisotropies are completely absent means that the model did not just accommodate cosmic time as a large-scale average---it empirically confirmed its necessity as a foundational prior to describe precise uniform expansion. As noted, this problem is epistemically equivalent to the expansion/existence, flatness, horizon, and energy budget problems: it is not that the theory is incapable of describing our observations, but rather that theoretically the observations we've made should not be allowed.

The observed exact isotropy of expansion does not simply rule out expansion anisotropies due to local density variations---it also fundamentally constrains the nature of structure formation itself. If local overdensities had any significant effect on the expansion rate of space, as standard GR-based reasoning suggests, then the expected perturbations to the Hubble flow should have been visible at orders of magnitude above the observed anisotropies. The complete absence of such an effect suggests that density variations do not influence expansion at all in the way the standard framework assumes. Instead, expansion must be an all-or-nothing phenomenon: structures either participate in uniform cosmic expansion or they do not. This naturally implies the existence of a threshold density---a critical value above which gravitationally bound structures cease to expand, while all structures below it participate fully in the uniform expansion of space.

Empirical evidence across multiple scales substantiates this threshold concept. The Local Group, which has a density surpassing the cosmological critical density by several times, remains fully gravitationally bound and does not participate in expansion. However, the Local Group–Virgo Cluster region, though still highly overdense, is already below this threshold, meaning Virgo is expanding away from us along with the Hubble flow. At even larger scales, Bo\"{o}tes Void galaxies, despite residing in one of the largest known underdensities, adhere strictly to the uniform Hubble flow rather than exhibiting residual perturbations from surrounding density variations \cite{Kirshner1981, Szomoru1996}. Studies have further identified bound objects within voids, indicating that even in such extreme underdense environments, localized regions exceeding the threshold can remain gravitationally self-contained \cite{Rieder2013, Alpaslan2014}. Finally, bound structures within Basin-of-Attraction (BoA) scale inhomogeneities continue evolving gravitationally despite their participation in cosmic expansion---suggesting that their expansion is dictated solely by the threshold condition rather than by density perturbations at any given scale \cite{Tully2014, Valade2024}.

These observations challenge the standard expectation that local density variations perturb the expansion rate, inevitably introducing anisotropies. Instead, the consistent uniformity of expansion observed across all sub-threshold structures implies that local inhomogeneities do not influence the overall expansion dynamics. The absence of any residual expansion perturbations at observable scales further reinforces the necessity of an absolute global expansion frame, confirming that the assumption of statistically uniform expansion was not just a convenient approximation but an exact empirical reality. The existence of a threshold density strengthens this conclusion, as it demonstrates that structure formation itself is constrained by this absolute expansion framework. Rather than treating uniform expansion as an assumption that holds on average, the data instead reveal it to be a fundamental geometric principle that governs cosmic evolution at all scales.

The cosmic time problem follows immediately from the isotropy problem and this threshold phenomenon. In order for expansion-induced redshifts to go undetected relative to the small-scale $3\times10^{-6}$ perturbations arising from the SW effect, our universe must be expanding uniformly outside of bound structures. Otherwise, we would require a fine-tuning mechanism to dampen the redshift anisotropy by more than $10^3$ so that it could go undetected. Alternatively, we might trust the evidence that our universe's expansion, which is itself an unexplained phenomenon, actually is uniform---as the standard model describes. In any event---whether extraordinarily fine-tuned or exact---this is a phenomenon we have empirically verified with several orders of magnitude better than expected precision. Our standard model, and its assumptions about the isotropy, homogeneity, and uniform expansion of space, has been precisely verified. And those assumptions are all predicated on the more basic assumption that there is a cosmic time. The empirical verification of precisely uniform expansion is, at its core, the empirical verification of a global cosmic time. If our universe did not have a well-defined cosmic time, the observed uniformity of expansion would be impossible. The precision with which we confirm uniform expansion demands that any alternative model lacking a well-defined cosmic time be ruled out at the same confidence level. Thus, what was originally a mere convention---the assumption that a large-scale average cosmic time exists---has now been elevated to a precise empirical necessity. The standard model has not just assumed cosmic time; it has inadvertently verified it with extraordinary precision.

The precise isotropy of the CMB monopole is not merely evidence of an average thermodynamic reference frame at recombination---it is a direct empirical constraint on the nature of cosmic expansion itself. The fact that this preferred state of rest has persisted unaltered for 13.8 billion years is not an assumption that GR predicts---it is an empirical fact that GR must now explain. Our \textit{hope} that at smaller scales matter should affect the expansion rate while the cosmological rate would describe only a large-scale approximation has been thoroughly falsified by the precision of the monopole we observe once the dipole, the Milky Way contamination, and all secondary gravitational effects have been accounted for. Any expansion-induced anisotropy that may exist is completely overwhelmed by these secondary effects at the $3\times10^{-6}$-level, even though a conservative estimate based on the observed density contrast predicts that they should be a thousand times more pronounced. The evidence is screaming at us through loud speakers, telling us that we live in an enduring, uniformly expanding, maximally symmetric well-defined three-dimensional universe in which time passes with a well-defined Lorentzian metric connection.

The original justification of the line-element whose properties---cosmic time, isotropy, homogeneity, and uniform expansion---we have now significantly constrained, is worth revisiting \cite{Robertson1933}: ``That we require a spatial distribution of matter in our highly idealized universe to be uniform implies the existence of a significant simultaneity, and would at first seem contrary to the postulates of the theory of relativity, according to which each observer refers the world to his own proper space and time. But this difficulty disappears on introducing, in accordance with the observed facts, the assumption that there exists in each region of cosmic space-time a \textit{mean} motion which represents the actual motions to within relatively small and unsystematic deviations.'' Now, nearly a century later, this tension Robertson noted but brushed off with the notion that the contradictory principles should only hold as a large-scale average in our local region of cosmic space-time, has been shown to have actually persisted with extraordinary precision, across all of observable space, throughout the entire history of the observable universe.

Against the standard operationalist interpretation of relativity the evidence stands: in reality there is an absolute frame of reference, an absolute cosmic time, and, outside of bound structures, an absolute, uniform expansion of our maximally-symmetric three-dimensional universe. 

The notion, based on an interpretation of the principle of relativity, that such a ``preferred'' frame of rest could never be observed, and that one's velocity relative to it could never be measured, is strongly contradicted by the empirical evidence. The standard model assumes that cosmic expansion should be approximately isotropic on large scales, yet the data reveal that isotropy is not just an approximation---it is an exact feature of our universe. This is why we can precisely measure our velocity relative to the absolute cosmic frame---$v=(369.82\pm0.11)$ km/s towards the constellation Crater \cite{Aghanim2020a}---through a $\delta T/T \approx 3 \times 10^{-3}$ dipole. However, in a universe where expansion anisotropies due to structure formation were present at the expected $10^{-2}$ level, such a clean dipole should have been obscured by large-scale redshift distortions that do not appear in the data. This is possible because physical processes in our expanding universe produced the CMB \textit{and then continually preserved} its initial isotropy through isotropic expansion over the past 13.8 billion years. These observations suggest that the standard operationalist interpretation of relativity, which denies the existence of a preferred cosmic time, may be incomplete. 

While GR as a mathematical framework is consistent with a preferred foliation of space-time---such as that described by the FLRW model---the broader relativistic paradigm has historically rejected such a notion. The extreme precision with which isotropic expansion has been verified now challenges that rejection, indicating that a well-defined global cosmic time is not merely a convenient convention but an empirical reality. If relativity is to remain the foundation of cosmology, it must be reconciled with this fundamental constraint. As with the expansion/existence problem, the horizon problem, the flatness problem, the energy budget problem, and the isotropy problem, the epistemic issue is that the absolute foliation that evidently exists is not \textit{supposed} to exist, and stands in direct conflict with the generality that forms the basis of the standard operationalist interpretation of GR. Thus, the very idea that relativity abhors a preferred cosmic time is not just unwarranted---it is directly contradicted by the empirical evidence.

The standard model does not merely lack a theoretical explanation for the observational constraints related to the six epistemically equivalent problems discussed here---it actively leads to the expectation that these properties \textit{should not} be observed without extreme fine-tuning. While inflation, dark energy, and other proposed modifications attempt to resolve individual problems within the standard framework, they do so by introducing additional assumptions rather than addressing the deeper conceptual issue: that GR, as currently formulated, lacks a natural mechanism to account for these properties. If the only way to make GR compatible with cosmological observations is through a growing number of fine-tuned, ad hoc fixes, then this suggests a more fundamental revision of the theory may be necessary.

For over a century, cosmology has been explored within a relativistic framework, treating standard operationalist GR as the foundation upon which our understanding of the universe is built. Yet, what we have found is that the data repeatedly reveal tensions between GR-based expectations and the observed large-scale structure of the cosmos. Each of the six epistemically equivalent problems identified here---whether related to expansion, isotropy, or cosmic time---points to a fundamental misalignment between theoretical assumptions and empirical constraints.

If the only way to make GR compatible with cosmological observations is through a growing number of fine-tuned, ad hoc fixes, then this is not merely a sign that the theory is incomplete---it is an indication that we have been looking at the problem backwards. In light of these tensions, it may finally be time to flip the script: rather than continuing to search for \textit{relativistic} solutions to \textit{cosmological} problems, treating the latter as an application of the former, we should instead consider what it would mean to develop a cosmological theory of relativity---one that takes the empirically-evident structure of our universe as primary and re-examines the theoretical framework accordingly.

\section{Cosmological Relativity}

Over the past century, numerous modifications to GR have been proposed in response to its known limitations, particularly in the cosmological regime. Alternative theories of gravity, such as scalar-tensor theories \cite{Brans1961}, $f(R)$ gravity \cite{Sotiriou2010}, and bimetric gravity \cite{Hassan2012}, introduce additional degrees of freedom to account for observed phenomena that GR alone struggles to explain. Meanwhile, attempts to address cosmic acceleration have led to proposals such as quintessence \cite{Ratra1988} and $f(T)$ gravity \cite{Cai2016}, while inhomogeneous models \cite{Buchert2000, Wiltshire2009, Seifert2025} have challenged the necessity of dark energy altogether. On smaller scales, Modified Newtonian Dynamics (MOND) \cite{Milgrom1983} and its relativistic extensions, such as TeVeS \cite{Bekenstein2004}, seek to explain galaxy rotation curves without invoking dark matter. Each of these approaches seeks to address empirical gaps within the standard model, yet none have succeeded in fully resolving the core tensions without introducing new theoretical difficulties.

However, the need for a conceptual revision may run even deeper than these approaches recognise. The issues discussed in the previous section---the isotropy and cosmic time problems among others---do not merely highlight empirical gaps in our current model but reveal a fundamental misalignment between GR's operationalist foundations and the cosmological structure of reality. Addressing these problems does not require introducing new dynamical fields or modifying the Einstein field equations. Instead, they may be resolved through a simple conceptual augmentation of GR that establishes a layered geometric structure to space-time---one that maintains GR's empirical successes while fundamentally redefining its physical meaning and reshaping its derived physical consequences---and resolves its cosmological tensions. This framework, which we refer to as \textit{Cosmological Relativity} (CR), builds on the insights of the ADM decomposition and ultimately draws an explicit distinction between noumena and phenomena, providing a minimal but necessary correction to GR.

\subsection{A layered geometric framework}

The layered geometric framework (LGF) of CR reinterprets space-time not as a fundamental entity, but as a derived structure that encodes the history of physical events occurring within the three-dimensional universe as cosmic time passes. It explicitly rejects the notion that space-time is the ``fabric of the cosmos,'' instead positing that the true fabric is the well-defined, uniformly expanding and enduring, three-dimensional universe. Rather than viewing space-time as an elastic, dynamical entity that governs motion, LGF posits that space itself is fundamental, and that its dynamically evolving curvature---determined by the distribution of matter---affects local perceptions of time's passage. 

The empirical constraints from the isotropy and cosmic time problems force a preferred foliation with absolute simultaneity, while the observed phenomena of gravitational waves, lensing, and time dilation demand a framework that allows real space to dynamically warp while preserving a maximally symmetric expanding background. This requires not only a preferred foliation but also diffeomorphic flexibility between real space and cosmological spatial slices. Together, these constraints uniquely determine the layered structure of CR as a necessary augmentation of GR.

In this framework, space-time remains a Lorentzian manifold \( (M, g_{\mu\nu}) \), while the physical universe is represented by a preferred foliation of space-like hypersurfaces \( \Sigma_t \), which evolve over cosmic time \( t \). These hypersurfaces are equipped with a Riemannian 3-metric \( h_{ij} \) that describes spatial distances at a given cosmic time. The lapse function \( N(x) \) and shift vector \( N^i(x) \) relate adjacent hypersurfaces, defining the temporal and spatial displacement between them. The mapping between the real three-dimensional universe and the ADM foliation must be diffeomorphic, ensuring that while \( h_{ij} \) is not necessarily identical to the physical spatial geometry, it remains smoothly mappable to it. 

A key consequence of LGF is that it maintains a Lorentzian metric connection between events that occur in space over time, ensuring that photon worldlines correspond to invariant null geodesics in the derived space-time structure. This metrical connection is what fundamentally distinguishes LGF from Newtonian absolute space while preserving relativistic causal constraints. LGF thereby retains the empirical success of general relativity while offering a revised geometric interpretation---one in which the universe is fundamentally three-dimensional, evolving over cosmic time, with space itself responding to the presence of matter and motion.

This distinction is crucial: it allows space-time to manifest differently at different scales, ensuring that the physics captured at one scale need not impose artificial constraints on another. The local curvature experienced in the vicinity of a massive body does not necessarily dictate the global structure of the expanding universe, nor does the large-scale homogeneity of cosmic expansion constrain the local space-time dynamics of gravitational collapse. By allowing distinct yet compatible layers, LGF provides the necessary flexibility to describe both local gravitational physics and global cosmology without forcing them into a single geometric framework that may be ill-suited to both. A direct consequence of this flexibility is the natural emergence of a threshold density that governs structure formation. In LGF, expansion is not modulated continuously by local inhomogeneities, as GR-based expectations would suggest, but rather follows a binary criterion: regions above the threshold density remain gravitationally bound, while regions below it participate fully in the uniform expansion of space. This aligns precisely with the empirical evidence across multiple scales, from the Local Group–Virgo transition to Bo\"{o}tes Void and BoA structures, reinforcing the idea that cosmic expansion operates within a sharply defined framework rather than through gradual dynamical modulation. Unlike standard GR-based interpretations, which struggle to reconcile large-scale expansion with the hierarchical formation of structure, LGF accommodates both within a single geometric description. The threshold density emerges naturally as a fundamental boundary condition of the universe's evolving spatial geometry, ensuring a consistent treatment of structure formation and cosmic flow.

Thus, LGF fundamentally refines the nature of simultaneity. Regardless of the coordinate frame used to describe space-time, the universe and the set of events that simultaneously occurred within it are described by a particular foliation, which may not be synchronous in any given frame of reference. This foliation must be consistently accounted for when interpreting dynamical sequences of spatially separated events across different space-time descriptions, as failing to do so may lead to incorrect conclusions about the temporal evolution of events in the universe. 

While an absolute foliation alone provides a useful conceptual shift, it does not by itself lead to significant new physical insights. The power of LGF arises not merely from the existence of a preferred foliation but from the recognition that the foliated spaces of cosmological space-time need not correspond exactly to the physical geometry of space itself---only that they must be diffeomorphic to it. In other words, local space-time and cosmological space-time can be geometrically distinct, yet remain mathematically related through smooth mappings that preserve their essential topological structure.

In essence, this layered approach restores a critical missing element in the standard operational interpretation of GR: the ability to separate phenomenological space-time, which describes what is observed in different regimes, from noumenal space, which encodes the deeper structure of reality. The failure to make this distinction in standard cosmology has led to an artificial unification of local and cosmological physics, where assumptions valid at one scale are unjustifiably extrapolated to another. By recognising that local and cosmological space-time may exhibit distinct geometric structures at different scales while remaining diffeomorphic, LGF provides the conceptual framework necessary to address the deep tensions between GR and cosmology. 

This framework is not merely an alternative formalism---it is an empirically necessary extension of GR. The assumptions of operational relativity (OR) fail to reconcile the significant empirical constraints on uniform cosmic time and expansion with local phenomena like gravitational waves and lenses, but CR, through its layered structure, provides the minimal mathematical augmentation required to resolve these contradictions while preserving all known empirical successes of GR.

This framework not only resolves the isotropy and cosmic time problems but also naturally incorporates the existence of a critical density threshold, explaining why large-scale structures exhibit precisely uniform expansion while bound structures evolve gravitationally without contributing to the Hubble flow. Rather than treating this threshold as an emergent feature of complex gravitational interactions, LGF recognises it as an intrinsic aspect of cosmic geometry---one that follows from the fundamental distinction between local and global space-time dynamics. This perspective ensures that cosmic expansion is neither a purely local phenomenon nor an arbitrary large-scale approximation but an absolute geometric principle governing the universe at all observable scales.

\subsection{Discussion: deconflating phenomena and noumena}\label{sec:noum}

When we look at the horizon, the Earth appears flat. Yet the changing altitude of stars as we move north or south, along with Earth's curved shadow during a lunar eclipse, reveals that Earth is actually a sphere. The stars, too, appear to move across the sky each night in perfect concentric arcs, leading people to assume for millennia that they physically orbited Earth. But Copernicus recognized that this motion was an illusion---the stars remain fixed while Earth itself rotates. The world does not \textit{feel} like it is moving---neither spinning on its axis nor racing through space---but Galileo explained why: inertial motion is imperceptible.

The Sun and planets, for their part, \textit{appear} to trace complex paths against the stars, sometimes slowing, stopping, and even reversing course. For over a thousand years, astronomers built intricate systems of epicycles, deferents, eccentrics and equants to model this motion as real. But Kepler revealed the deeper truth: planets do not move in loops across the sky. The retrograde motion we see is entirely an illusion, caused by our own motion around the Sun.

This pattern repeats throughout physics: we begin with models that match observations directly, but when those models fail, we are forced to recognize that reality is fundamentally different from how it appears.

The Earth \textit{appears} flat, but it is not. The spherical model not only accounts for the horizon's apparent flatness but explains a range of other observations, from star positions to images we've directly taken from space. The stars \textit{appear} to orbit Earth, but their motion is fully explained by Earth's rotation. Though we \textit{feel} at rest, we are not---and our understanding of inertia tells us exactly why. The Sun and planets \textit{appear} to drift through the stars, with the planets sometimes reversing course---but this, too, is an illusion caused by our own motion around the Sun.

None of these discoveries were obvious. Each revealed that reality is fundamentally different from how it appears. Most importantly, each required a formal distinction between \textit{phenomena}---what we directly observe---and \textit{noumena}---the underlying reality that gives rise to those observations.

Relativity theory has followed a remarkably similar path. Einstein's operational definition of simultaneity equated it directly to frame-dependent synchronicity \cite{Einstein1905a,Einstein1952}. This definition---along with the assertion that if a preferred frame \textit{did} exist, it could never be observed---became sacrosanct. And so, while we have been modelling our universe with a preferred frame for a century, and refining our measured velocity through it for decades, we have continued to insist that this frame is merely a convenient choice---a large-scale average---rather than a fundamental aspect of reality.

Modern relativistic cosmology has inherited this assumption uncritically. Just as early astronomers assumed that planetary retrograde motion was a real phenomenon, OR assumes that space-time itself---the full set of events that ever happen \textit{in} reality---\textit{is} reality itself---the ``fabric of the cosmos''---rather than a description of a deeper evolving structure---the three-dimensional universe we model, which expands and evolves dynamically over time. The very epistemic principle that was invalidated in the Scientific Revolution---that things appear the way they do because that is how they are---has been smuggled back into physics under a different name: operationalism. Instead of questioning this assumption when confronted with contradictions, the standard model simply absorbed those contradictions by constructing ad hoc fixes to explain them.

But this is precisely where the FLRW model runs us afoul. In formulating the standard model, we defined a preferred foliation based on a large-scale average cosmic time---an absolute simultaneity---but then we further imposed the operationally relativistic assumption that this cosmic time must be synchronous in the cosmic rest frame. This constraint is neither required nor fundamental; it is an artificial imposition of an operationalist principle onto a general relativistic universe. And the deeper issue is that by enforcing synchrony in cosmic time, the standard model makes the same naive and superficial heuristic assumption that proved so problematic in the past: as the model describes, the observed expansion rate \textit{is} the fundamental expansion rate, the inferred curvature \textit{is} the actual curvature, and the measured density parameters \textit{are} the true fundamental densities of the universe. Just as planetary epicycles were effective empirical parameters that described apparent retrograde motion without revealing the deeper heliocentric structure, the standard model's cosmological parameters may likewise be effective phenomenological quantities rather than fundamental noumenal properties. The model's implicit conflation of these categories obscures the possibility that the parameters we fit to observations are merely empirical values that emerge from a deeper theoretical structure we have yet to recognise.

But this conflation is not just a conceptual oversight---it is the fundamental structural weakness of the relativistic framework as currently formulated. GR attempts to describe everything at once: local space-time warping, global cosmic evolution, and even singular structures like black holes, all within a single continuous manifold. But in doing so, it mashes together conceptually distinct elements: the evolving cosmological structure of the universe, the local warping of space-time, and the causal constraints on lightcones, forcing them into a singular, unexamined model. This lack of separation leads to incoherence: we simultaneously claim that space-time is locally Minkowskian while insisting on a global FLRW structure, treating black holes as fully formed in an FLRW universe despite the fact that external observers can never see collapse complete.

This structural failure is not a problem with GR itself---it is a problem with how GR has been applied to cosmology. The missing element is not new physics, but a diffeomorphic flexibility that allows for distinct but internally coherent representations of different layers of reality. CR rectifies this, not by rejecting GR, but by formally distinguishing between these layers: the global cosmic evolution is governed by an absolute cosmic time, while local space-time warping remains diffeomorphic to that foliation. Rather than forcing an operationalist, synchronicity-based interpretation onto the entire universe, CR introduces the necessary separation of structure that makes sense of what we already model tacitly, but without consistency.

This is the fatal flaw of the operationalist model: it conflated a phenomenological description with fundamental reality. The empirical confirmation of precisely uniform expansion is not just a consistency check on the standard model---it is the direct verification of an absolute cosmic time.

And once absolute cosmic time is accepted as a physical reality, the synchronicity assumption is not just unnecessary---it is an artificial imposition that has led us astray, in the same way we've found ourselves led astray in the past. The assumption that simultaneity must be synchronous in the cosmic frame is as naive as the geocentric assumption that planetary retrograde motion is real. It is an archaic vestige of an earlier, incomplete understanding, and its continued presence in the standard model reflects not a deep truth about the universe, but a historical artifact of an outdated operationalist framework. CR, which explicitly commits to the absolute cosmic time implied by cosmological evidence, formally rendering the synchronicity assumption arbitrary and unnecessary, is not a rival to GR---it is the necessary refinement that allows it to make sense of its own empirical successes. It does not replace relativity; it completes it.

\section{Applications of cosmological relativity}\label{sec:app}

While CR and GR share the same fundamental field equations, their mathematical frameworks are formally distinct in their treatment of space-time structure. CR explicitly includes a preferred foliation and relaxes the assumption that space-time slices must correspond directly to physical three-dimensional space at any cosmic instant, requiring only a diffeomorphic relation. These structural differences lead to distinct physical interpretations and applications.

For example, if space is flat in one reference frame then CR says that the most objective description of space that can be given from the perspective of an observer at rest in that frame is that space really is flat. CR rejects this ontological commitment. Similarly, if two phenomena occur at the same time in a given reference frame, CR asserts that the most objective possible description is that these events truly were simultaneous---while acknowledging that this description is frame-dependent. CR, on the other hand, formally commits to a different ontology in which simultaneity is objectively defined globally, and is coordinate system-independent.

These formal distinctions between the two frameworks lead to significant differences in the way physical phenomena are described. The following four examples illustrate these differences, showing how CR resolves longstanding conceptual failures of CR while preserving its phenomenological success.

\subsection{Euclidean cosmology}\label{sec:EC}

Consider a large trampoline, with bowling balls of varying densities rolling around on it, growing and shrinking and merging together. The denser these balls become, the more the trampoline warps under their concentrated weight. In this intuitive picture, each moment of the trampoline's existence is well defined, and as long as one of the balls does not shrink to infinite density and punch a hole in it, at each moment it is always diffeomorphic to Euclidean space. 

We will return to the problem of these balls' potential to shrink to inifinite density in \S~\ref{sec:BH}. For now, let us assume this cannot physically happen within finite cosmic time according to CR.

In this thought experiment, nothing has yet been said about the measure of the trampoline's existence---i.e., the metrical passage of time---nor how that connects to the propagation of information or the rates at which the balls and ripples can move. To establish this structure, we impose a single physical constraint: ripples on the trampoline, which carry information, always propagate at a fixed speed in all directions, independent of the motion of their source. Meanwhile, the massive balls move more slowly than the ripples, defining a clear separation between massless and massive objects. This is equivalent to enforcing the light postulate of special relativity, ensuring that all observers agree on the limiting speed of causal influence.

From these constraints alone, the full Minkowski space-time structure follows naturally as the cosmological space-time metric. Since the instantaneous shape of the trampoline is always diffeomorphic to Euclidean space, and since information must propagate along null lines, the underlying space-time metric must be Lorentzian, with a well-defined causal structure. The only metric compatible with this setup---where space is globally Euclidean at each moment but evolves dynamically, with an invariant speed limit---is Minkowski space-time.

This conclusion remains unchanged when replacing the two-dimensional trampoline surface with an infinite three-dimensional space in which massive objects move and warp their surroundings while ripples propagate at finite speed. The cosmological space-time metric of special relativity naturally emerges as an infinite Minkowski metric, even though the real three-dimensional universe is dynamically evolving and locally warping due to matter densities. In other words, Minkowski space-time is not the physical universe itself, but rather the cosmological background within which the real three-dimensional universe exists and evolves.

In this dynamical Euclidean cosmology, where the large-scale space-time structure remains precisely Minkowski, gravitational waves propagate as real, evolving distortions of space rather than abstract perturbations of a four-dimensional metric. The trampoline analogy clarifies this: just as ripples on a two-dimensional surface propagate within the surface itself, gravitational waves in CR exist within the evolving three-dimensional universe, moving through space over time rather than oscillating as features of a static space-time fabric. Unlike CR, which treats space-time as a fixed four-dimensional structure where gravitational waves are merely variations in the metric, CR ensures that these waves exist within a real, Lorentzian medium---space itself---eliminating the ambiguity of “what is waving?” and providing a clearer foundation for gravitational wave physics. This interpretation naturally resolves the problem of gravitational wave energy transfer, as energy is carried by real spatial distortions rather than coordinate-dependent fluctuations in a metric. By embedding gravitational waves within an evolving three-dimensional universe, CR restores a physical, intuitive picture of their propagation while remaining diffeomorphic to GR.

While gravitational waves describe the fabric of space itself evolving over time, other physical wave phenomena propagate through space rather than distorting it. These, too, naturally fit within the CR framework. Electromagnetism, like gravitational waves, is inherently Lorentzian and formulated in a 3+1 structure, but CR further clarifies its role by embedding electric and magnetic fields within an evolving three-dimensional universe rather than treating them as tensor components of a static four-dimensional metric. Meanwhile, fluid dynamics, plasma physics, and lattice gauge theories---which all rely on an evolving 3+1 formulation---struggle to fit naturally into CR's rigid four-dimensional structure. In CR, these theories find a more consistent setting, where their fundamental equations remain valid while avoiding the artificial constraints of embedding them in a global space-time manifold. Similarly, quantum mechanics, which fundamentally assumes a 3+1 formulation for wavefunction evolution, is more naturally interpreted within a dynamically evolving three-dimensional space rather than a frozen four-dimensional block universe. These structural advantages suggest that CR may provide a more physically coherent foundation for multiple domains of physics that have long relied on heuristic 3+1 formulations, yet have lacked a fully consistent relativistic setting.

This example formalizes the distinction between OR and CR, reinforcing that the two frameworks are not merely different perspectives but structurally distinct theories of space-time. Unlike OR, which assumes space-time slices correspond directly to physical three-dimensional space, CR introduces an absolute simultaneity while allowing the physical universe to be only diffeomorphic to the foliated spatial slices of the cosmological space-time metric, rather than strictly identical to them. This distinction provides greater descriptive flexibility without violating diffeomorphism invariance, resolving the inconsistencies that arise from OR's rigid geometric constraints which deliberately conflate phenomena and noumena at a superficial level. More importantly, beyond its empirical motivation from cosmology, CR demonstrates its greater structural power by accommodating a broader range of physical wave phenomena within a unified relativistic framework. The fact that CR emerges naturally from first principles while preserving all phenomenological successes of OR and extending its applicability to new domains of physics underscores its necessity---not as an alternative theory, but as the more complete and physically coherent foundation for relativity.

\subsection{Time Travel}

CR eliminates time travel at the fundamental level by rejecting the notion that space-time is reality. In CR, time travel paradoxes arise because simultaneity is frame-dependent---the notion of what events are happening ``now'' shifts depending on motion. This allows for theoretical causal loops when superluminal travel is introduced or when general relativity's global space-time admits closed timelike curves. However, CR renders these scenarios physically impossible by establishing absolute simultaneity and eliminating the block universe interpretation.  

In CR, the three-dimensional universe exists in an absolute cosmic time, and all physical evolution happens within this evolving three-dimensional structure. There is no pre-existing four-dimensional space-time manifold that allows causal loops to be embedded---space evolves forward over time, and any paradoxes arising from relativity of simultaneity disappear at the ontological level.  

Because simultaneity in CR is globally well-defined and absolute, the conditions necessary to construct a time travel paradox simply do not exist:  
\begin{enumerate}
    \item \textbf{Superluminal travel cannot lead to backwards-in-time motion}---an event that happens at cosmic time $t$ is objectively in the past or future of another event, independent of any observer's motion.  
    \item \textbf{Closed timelike curves are physically meaningless}---since CR does not describe the universe as a fixed four-dimensional manifold, no self-consistent looping trajectories can be embedded in space-time. Any trajectory must be a continuous evolution within the three-dimensional universe layer, eliminating any possibility of time loops.  
    \item \textbf{Time travel paradoxes rely on the operationalist conflation of phenomena and noumena}---the assumption that coordinate-dependent effects like relativity of simultaneity define physical reality. But in CR, simultaneity is objectively defined at the cosmological level, so paradoxes like the grandfather paradox never arise.  
\end{enumerate}  

In short, CR does not ``resolve'' time travel paradoxes---it makes them physically impossible. The entire framework that allows for time travel in CR does not exist in CR because the notion that space-time is reality has been rejected, simultaneity is absolute, and space evolves forward in a well-defined cosmic time. Time travel is not just prohibited---it is ontologically meaningless in a theory that distinguishes the mathematical description of space-time from the real, evolving structure of the universe. 

\subsection{Gravitational collapse}\label{sec:BH}

\begin{figure*}[t]
\centering
\includegraphics[width=3.3in]{EF_v-r.png}\includegraphics[width=3.3in]{EF_Sch_t.png}\\
\hspace{0.1in}(a)\hspace{3.05in}(b)
\caption{Space-time diagrams depicting the causal structure of the Schwarzschild geometry in the Eddington-Finkelstein reference frame. (a) Constant $v-r$ slices as the chosen foliation, where outgoing (red) and infalling (blue) null geodesics are depicted alongside equal-interval spatial slices (black). (b) Constant $t$ slices defining the foliation, where all infalling matter asymptotically approaches the event horizon at $t\rightarrow\infty$.}
\label{fig:BH_EF}
\end{figure*}

The standard picture of spherically symmetric gravitational collapse to a black hole is essentially motivated with reference to Figure \ref{fig:BH_EF}(a), on which timelike worldlines of infalling particles and outside observers alike can be drawn within the bounds of the causal framework. Technically, the basic mathematical framework of GR does not specify which points along two separate worldlines should ever be understood to occur at the same time, but the idea that an external observer \textit{remains} outside \textit{after} an infalling particle has crossed the horizon or reached the singularity at $r=0$ requires such an equivocation to be made. So while GR technically lacks this infrastructure, OR-based intuition tends to fill the gap so that the standard picture of gravitational collapse uncritically assumes simultaneity of synchronous events in such reference frames. Thus, black lines of constant $v-r$ are drawn explicitly on Figure \ref{fig:BH_EF}a to complete the framework leading to the standard concepual model of black holes, providing the background infrastructure on which one might draw the collapsing timelike surface of a star, or an outside space ship and an infalling space ship. 

When Penrose described the physical process of gravitational collapse alongside his formulation of the cosmic censorship hypothesis \cite{Penrose1969,penrose2002}, he implicitly made this OR-based equivocation, stating that ``No matter how long the external observer waits, he can always (in principle) still see the surface of the star as it \textit{was} just before it plunged through the Schwarzschild radius.'' This statement implicitly assumes an interpretation of simultaneity according to which the surface of the star crosses the horizon in finite time, even though this is not globally true, according to which the external observer \textit{remains} outside \textit{after} the infalling star has \textit{already} passed through the horizon. In effect, Penrose's description frames black holes as space-time regions in which collapse has already completed in a global sense, reinforcing the standard but unexamined intuition that the horizon represents an objectively real boundary that has formed in finite time.

However, this conclusion does not follow logically from GR itself. Different spatial foliations describe the infalling star differently; in particular, foliations such as Schwarzschild $t$-slicing describe the horizon as \textit{never} forming in finite time (see Figure~\ref{fig:BH_EF}b. Since GR \textit{does not} prescribe a unique global simultaneity, no single slicing can claim privileged status in determining what ``has already happened'' in an objective sense. For the standard black hole picture to be valid---to be an inescapable logical consequence---horizon formation in finite time must be true in \textit{every} possible foliation---but it is not. The interpretation that collapse completes in a finite, well-defined global time is not a prediction of GR, but rather an unexamined consequence of an unjustified OR-based equivocation.

To understand how spherically symmetric gravitational collapse must be handled in CR, we begin with the framework established in \S~\ref{sec:EC}: a Euclidean cosmological space-time, in which the universe $\Sigma_t$ may be locally warped in the presence of matter while remaining diffeomorphic to three-dimensional Euclidean space. Consider a spherically symmetric mass at rest in this space. The corresponding space-time line element in the vicinity of this mass must take the general form
\begin{equation} 
    ds^2=-A(r,t)dt^2+B(r,t)dr^2+r^2d\Omega^2, 
    \label{eq:JB1} 
\end{equation}
where $t$ is the time coordinate of the cosmological foliation, which we define to be identical to the coordinate time in this local description for a mass at rest, ensuring a direct mapping between the global cosmological slicing and the local metric. The coordinate $r$ is the radial coordinate in the exterior region of the mass, and $d\Omega^2$ represents the line element of the unit two-sphere, encoding the spherical symmetry of the system. To determine the specific form of this metric outside the mass, we solve the vacuum Einstein equations in this spherically symmetric setting. The result is the Jebsen-Birkhoff theorem \cite{Jebsen1921,Jebsen2005,Birkhoff1923}, which uniquely determines the local space-time geometry as the Schwarzschild solution in stationary coordinates:
\begin{equation} 
    ds^2=-\left(1-\frac{r_{\mathrm{Sch}}}{r}\right)dt^2+\frac{1}{1-\frac{r_{\mathrm{Sch}}}{r}}dr^2+r^2d\Omega^2, 
    \label{eq:Sch} 
\end{equation}
where $r_{\mathrm{Sch}} = \frac{2GM}{c^2}$ defines the Schwarzschild radius, which corresponds to a coordinate singularity (the so-called ``event horizon''), dependent on the mass $M$ of the central body, the gravitational constant $G$, and the speed of light $c$.

In the Schwarzschild solution, the coordinate time $t$ is defined in accordance with the global cosmological foliation, ensuring that all spatial slices remain diffeomorphic to three-dimensional Euclidean space. A fundamental feature of this slicing is that infalling material asymptotically approaches the Schwarzschild horizon at $r = r_{\mathrm{Sch}}$ but never reaches it within finite $t$. Since GR does not prescribe a unique simultaneity structure, an observer moving at constant velocity relative to the mass might attempt to redefine simultaneity using their own Lorentz-boosted coordinates. However, such a transformation cannot alter the fundamental asymptotic nature of the Schwarzschild slicing.

This follows from the fact that Lorentz transformations in flat space-time mix time and space coordinates, but they do not change causal structure in curved space-times where the event horizon remains an invariant feature of the geometry. Specifically:
\begin{enumerate}
    \item A Lorentz boost mixes Schwarzschild $t$ and $r$ coordinates but does not alter the asymptotic behavior of infalling geodesics. Even in a boosted frame, an object freely falling toward the horizon follows a worldline that asymptotically approaches $r_{\mathrm{Sch}}$ without ever crossing it in finite coordinate time.
    \item The global causal structure of Schwarzschild space-time is invariant under boosts. The Schwarzschild $t$-slicing defines a family of spatial hypersurfaces that extend arbitrarily close to $r_{\mathrm{Sch}}$ but never cross it, ensuring that all events at the horizon or below remain outside the coordinate domain of these slices. Since the Schwarzschild metric is static in these coordinates, applying a Lorentz boost to any observer outside the black hole cannot change this global slicing to one in which horizon crossing occurs at finite coordinate time.
    \item Even highly boosted observers cannot ``tilt'' space-time slices enough to force a crossing. The Schwarzschild slicing remains an asymptotic limit, and no constant-velocity transformation is sufficient to alter the foliation structure in such a way that infalling matter reaches $r_{\mathrm{Sch}}$ in finite time across all frames.
\end{enumerate}

Thus, the Schwarzschild foliation retains its fundamental structure regardless of any boosts, ensuring that the basic framework shown in Figure~\ref{fig:BH_EF}b essentially captures the cosmological dynamical evolution of motion through space surrounding a spherical mass. Accordingly, while the surface of a collapsing star will reach $r_{\mathrm{Sch}}$ in finite proper time, this interval requires infinite cosmic time. Furthermore, the entire collapsing star, as well as every other particle that ever falls towards the black hole, asymptotically approaches 
$r_{\mathrm{Sch}}$ as $t\rightarrow\infty$, meaning that in cosmic time, all infalling matter remains outside the event horizon for the entire lifetime of the universe.

According to CR, the standard conceptual model of gravitational collapse to a black hole in GR, represented by the framework in Figure~\ref{fig:BH_EF}a, should be replaced by the model derived from Figure~\ref{fig:BH_EF}b, according to which black holes in \textit{our} universe are not objects that have already formed, but are instead regions of asymptotic collapse, where infalling matter continuously approaches---but never reaches---the event horizon in finite cosmic time.

Since the Schwarzschild solution uniquely describes the vacuum exterior of a spherically symmetric mass, the surrounding space-time remains Schwarzschild throughout the collapse process, with the Schwarzschild radius $r_{\mathrm{Sch}}$increasing as additional mass accretes. However, in CR, the surface of the collapsing object itself asymptotically approaches but never reaches $r_{\mathrm{Sch}}$ in finite cosmic time. This means that, from a global perspective, black holes in \textit{our} universe are never fully formed objects but are instead asymptotically collapsing regions of space-time. While this interpretation is mathematically diffeomorphic to GR, it carries potential observational consequences: in particular, gravitational wave signals from black hole mergers may differ from standard expectations, since merging black holes are never fully-formed static objects but instead dynamically evolving regions of asymptotic collapse. This distinction could lead to subtle deviations in ringdown waveforms or inspiral dynamics compared to models that assume pre-existing event horizons, similar to known differences in gravitational wave signatures from neutron star versus black hole mergers. Future high-precision gravitational wave observatories may provide a means to empirically test this aspect of CR.

\subsection{dS/SdS cosmology}

The examples so far have assumed a cosmological space-time in which the three-dimensional universe is diffeomorphic to static Euclidean space. However, observational evidence indicates that our universe is not static but uniformly expanding with a well-defined cosmic time. To explore this, we adopt de~Sitter (dS) space as a basis for cosmological space-time.

4D dS space can be described as a hyperboloid of one sheet embedded in 5D Minkowski space, where the hyperboloid's axis is aligned with the time coordinate of the embedding space. This surface is doubly ruled, meaning it is covered by two complementary families of flat null geodesics. The global causal structure of dS space follows from this embedding, but its physical interpretation depends on the foliation imposed upon it.

A natural foliation, motivated by the hyperboloid's construction, is to align time slices with the embedding coordinate $T$, yielding a space-time of expanding and contracting 3-spheres:
\begin{equation}
    ds^2=-dT^2+\frac{3}{\Lambda}\cosh^2\left(\sqrt{\frac{\Lambda}{3}}T\right)\left(\frac{dR^2}{1-R^2}+R^2d\Omega^2\right).\label{eq:dS_hyp}
\end{equation}
For $T>0$, this metric meets the cosmological requirements of a maximally symmetric, uniformly expanding space. However, this formulation is problematic in two respects:
\begin{enumerate}
    \item The detailed expansion rate does not match observations.
    \item This space is synchronous in the cosmic rest frame, suffering the same fundamental shortcomings as FLRW cosmology discussed in \S~\ref{sec:noum}.
\end{enumerate}

Rather than assuming the standard causal structure of Eq.~(\ref{eq:dS_hyp}), we introduce a redefinition of the causal structure while preserving the same cosmic foliation. Specifically, we invert the roles of massive and massless worldlines:
\begin{itemize}
    \item The new fundamental cosmological rest frame is defined by the comoving bundle of null geodesics moving in one direction of the 3-spherical universe.
    \item The geodesics that remain at rest in Eq.~(\ref{eq:dS_hyp}) are reinterpreted as the null worldlines of photons.
\end{itemize}
This transformation does not alter the foliation of cosmic time or the fundamental symmetries of space, but it modifies the interpretation of motion within this space, fundamentally altering the metric structure (see Figure~\ref{fig:dS_SdS}).

\begin{figure*}[t]
\centering
\includegraphics[width=6.5in]{dS-SdS.png}
\caption{A single fundamental worldline within a cosmological model constructed from the flat null lines of de~Sitter space, whose motion is unaffected by the curvature of space-time. The causal structure of this model is redefined so that the bundle of lines parallel to the black worldline form timelike worldlines, the bundle of orthogonal lines shown here in red describe flat spacelike slices, and the blue worldlines that are at rest in the enduring 3-sphere become null lines. An inset shows the same set of lines displayed in the proper ($\chi,\tau$)-frame, where the `throat' of de~Sitter space is the singularity $\tau+\chi=0$. In this cosmological model, spherical spatial slices of de~Sitter space, taken to describe the universe at any cosmic time, are described by diagonal lines $\tau+\chi=$constant.}
\label{fig:dS_SdS}
\end{figure*}

To construct a space-time that respects both this new causal structure and the cosmic foliation, we introduce a key change: we treat the radius of the expanding 3-sphere as a timelike coordinate $r$. Since the cosmic 3-sphere is always orthogonal to this timelike direction, the spatial structure in the two dimensions where fundamental observers are at rest should remain unchanged. A general line-element satisfying these conditions can be written as:
\begin{equation}
    ds^2=-B(r,t)dr^2+A(r,t)dt^2+r^2d\Omega^2,\label{eq:JB}
\end{equation}
where $d\Omega^2$ describes the two spherically symmetric spatial dimensions that contract and expand with $r$, the universe's radius and timelike coordinate.

This redefined structure must still satisfy the vacuum Einstein equations with a positive cosmological constant, as it fundamentally describes a variation of dS space. Solving these equations constrains the functions $B(r,t)$ and $A(r,t)$, yielding the Schwarzschild-de Sitter (SdS) metric:
\begin{equation}
    \mathrm{d}s^2=-\frac{r}{\frac{\Lambda}{3}r^3+\frac{2GM}{c^2}-r}\mathrm{d}r^2+\frac{\frac{\Lambda}{3}r^3+\frac{2GM}{c^2}-r}{r}\mathrm{d}t^2+r^2\mathrm{d}\Omega^2,\label{eq:SdS_statical}
\end{equation}
which is a valid cosmological model for $r>0$ as timelike when the inequality
\begin{equation}
    \frac{\Lambda G^2 M^2}{c^4} > \frac{1}{9}
\end{equation}
is satisfied. This constraint imposes a minimum total mass necessary for the SdS solution to describe a cosmological model, requiring:
\begin{equation}
M > \frac{c^2}{3\sqrt{\Lambda G}} \approx 4.3 \times 10^{52}~\text{kg}.
\end{equation}
By comparison, the calculated baryonic mass of the observable universe, $M_{b,\text{horizon},0} = 1.5 \times 10^{53}~\text{kg}$, and the total mass including dark matter, $M_{m,\text{horizon},0} = 9.6 \times 10^{53}~\text{kg}$, both exceed this threshold. This confirms that the SdS model is consistent with known mass-energy constraints, providing a novel geometric perspective on the total mass distribution of the universe.

The coordinate singularity at $r=0$ in Eq.~(\ref{eq:SdS_statical}) does not correspond to a physical singularity but arises as a consequence of the causal redefinition. Unlike curvature singularities, such as the Schwarzschild singularity at $r=0$, this feature is an artifact of the coordinate system. The original dS space-time contains no singularities, and the introduction of $r$ as a timelike coordinate in SdS merely relocates the degeneracy of the coordinate system. This is analogous to how Schwarzschild coordinates become degenerate at the Schwarzschild horizon, even though the space-time is regular there in a more appropriate coordinate system.

This implies that the singularity does not represent a breakdown of the theory but rather marks the limit of the new coordinate system. Just as Schwarzschild coordinates break down at $r_{\mathrm{Sch}}$ but are extended by Eddington-Finkelstein coordinates, the coordinate singularity at $r=0$ could, in principle, be reformulated in an alternative coordinate system. However, because the foliation of cosmic time remains well-defined in the SdS formulation, no additional extension is necessary for the cosmological interpretation.

A key question that arises in this cosmological model is how an anisotropic space-time geometry can still lead to an isotropic observational universe. The resolution lies in the fact that the fundamental observers are comoving within a maximally symmetric space and that the speed of light is invariant.

First, while the slices of constant cosmic time $r$ are not isotropic and homogeneous in this frame, this results from the nontrivial evolution of fundamental observers along their geodesics rather than from an intrinsic inhomogeneity of space. By construction, space itself remains maximally symmetric as an expanding 3-sphere. The anisotropy is therefore a feature of the coordinate description rather than of the fundamental spatial structure.

Second, the comoving fundamental observers do not experience relative motion apart from cosmic expansion. Because light propagates along null geodesics, its speed must remain the same in all directions relative to these observers. This follows directly from the invariance of the speed of light: in any local frame, photons travel at $c$ in all directions, just as in special relativity. Consequently, fundamental observers necessarily perceive an isotropic universe, as all incoming light reaches them in an isotropic manner.

To understand the observational implications of this cosmological model, we examine the proper reference frame of fundamental observers. These observers correspond to the bundle of geodesics that remain at rest with respect to the vanishing of the gravitational potential in the Schwarzschild-de Sitter (SdS) metric. Their motion through $t$ over the course of $r$ is purely a consequence of the underlying geometry. The line-element in this coordinate frame is \cite{Janzen2015}:
\begin{equation}
    ds^2=-d\tau^2+(\partial_{\chi}r)^2d\chi^2+r^2d\Omega^2,\label{eq:SdS_flat}
\end{equation}
where
\begin{equation}
    r(\tau,\chi)=\left(\frac{6GM}{\Lambda c^2}\right)^{1/3}\sinh^{2/3}\left(\frac{3}{2}\sqrt{\frac{\Lambda}{3}}[\tau+\chi]\right),\label{eq:r_SdS}
\end{equation}
defined on the interval $\tau+\chi>0$.

This line-element exhibits several striking features. First, $r$ does not evolve with the proper time $\tau$ of any particular fundamental worldline, but with the parameter $c\tilde{\tau}\equiv\tau+\chi$, which is tilted with respect to the fundamental rest frame. That is, \textit{this universe is not synchronous in the fundamental rest frame}---a direct consequence of the causal redefinition.

Furthermore, while slices of constant $\tau$ are Euclidean (since $\partial_{\chi}rd\chi=dr$), these are not cosmological spatial slices. The lines perpendicular to the $\tau$-direction correspond to the complementary bundle of null geodesics, which are perpendicular at the ``throat'' of the de~Sitter hyperboloid to the bundle of worldlines defining fundamental observers (see Figure~\ref{fig:dS_SdS}). These orthogonal lines form a complementary reference frame, extending from a finite distance at $\chi=-\tau$ to $\chi=+\infty$, with synchronous events spanning cosmic time from the singularity at $r=0$ out to $r\rightarrow\infty$. 

A second remarkable feature of this line-element is that $r(\tilde{\tau})$ is \textit{identical} to the scale factor of the flat $\Lambda$CDM model:
\begin{equation} 
    \frac{a(\tilde{\tau})}{a_0}=\left(\frac{3H_0^2\Omega_{m,0}}{\Lambda c^2}\right)^{1/3}\sinh^{2/3}\left(\frac{3}{2}\sqrt{\frac{\Lambda c^2}{3}}\tilde{\tau}\right).\label{eq:scalefac} 
\end{equation} 
This is an astonishing result: the expansion of this universe, when described in the proper rest frame of either bundle of null geodesics that generate the full de~Sitter geometry, follows \textit{exactly} the phenomenological expansion rate measured in \textit{our} universe. That is, an observer in this universe, looking out at galaxies and at the relic afterglow of an early hot, dense state, would measure their universe to have precisely the same expansion rate as ours.

This identity is more than just a mathematical coincidence: it suggests that the standard FLRW phenomenology, typically assumed to require a synchronous space-time structure, instead emerges naturally from the causal structure redefinition applied within this cosmological model. Rather than requiring matter-energy content to dictate the universe's evolution, the SdS framework inherently enforces a uniform expansion history, with the same observational consequences as a flat $\Lambda$CDM universe. This implies that the large-scale expansion of the universe is not an effect of its contents, but a geometric property of space-time itself.

\section{Conclusion}

Einstein was an operationalist. His success came from formulating physical theories based solely on observable quantities, allowing him to develop special and general relativity without invoking absolute structures. This approach was a powerful tool, but it also led to a crucial misinterpretation: treating the full set of observed events---space-time itself---as fundamental reality, rather than as a phenomenological description of an evolving universe. He explicitly defended a static space-time picture as ``more convenient and more objective,'' arguing that relativity compels us to favour this view over a traditional dynamical picture of motion through space \cite{Einstein1938}. This preference, while mathematically valid, was a choice of interpretation, not an inescapable physical truth. Now, however, the large-scale observational constraints that reveal its limitations allow us to correct this misinterpretation and refine our understanding of relativity.

In retrospect, when satellites were first launched to measure CMB temperature anisotropies, we should have expected to see randomly oriented anisotropies at the 1\% level or higher. This would have been a striking confirmation of Einstein's theory, aligning with the fundamental prediction of FLRW cosmology: that while cosmic time and uniform expansion serve as reasonable large-scale approximations, the inhomogeneous distribution of matter should have naturally induced expansion anisotropies, imprinting subtle variations in the overall redshift of the CMB.

But this was not predicted---and it is not what was observed.

Instead, precision CMB measurements revealed that no such effect exists. The observed isotropy of the CMB monopole is not merely a rough approximation---it is extraordinarily precise, persisting at least three orders of magnitude beyond theoretical expectation. This is not just an unexplained feature of cosmology; it is a \textit{direct empirical contradiction} that exposes a fundamental failure of GR-based reasoning.

This contradiction reveals an even deeper problem: the way relativists interpret space-time itself. In the strictest reading of operationalist relativity, all physical truth is confined to an observer's causal past; any description of reality beyond what is observable is formally meaningless. Yet, in practice, relativists violate this principle constantly. The very idea that a black hole exists in the present---despite the fact that its event horizon can never be observed to form within the causal past of any external observer---relies on privileging one simultaneity slicing over others. Likewise, cosmological models that attempt to reconcile inhomogeneous structure with an evolving space---such as timescape cosmology---implicitly patch together a preferred notion of ``now'' while denying that they have done so. This tacit assumption is universal: every practical application of relativity requires an inferred ``now,'' yet standard relativity lacks the structure to define it coherently. CR resolves this by making explicit what is already implicitly required---a global cosmic foliation---while maintaining full empirical consistency with relativity's tested predictions.

In this sense, CR is not just an alternative framework---it is the objective, logically consistent formalization of what has always been done tacitly, but without a coherent basis. Every time relativists discuss black holes as already formed objects, or describe the evolution of the universe using a patchwork of different proper times stitched together, they are implicitly inserting an inferred global ``now'' while simultaneously denying the legitimacy of doing so. This inconsistency has persisted not because it is justified, but because it was never properly examined. CR does nothing more than formalize this necessary structure in an explicit and principled way---ensuring that relativity is applied coherently across all scales rather than selectively and surreptitiously.

The accidental verification of precisely uniform expansion, against general relativistic expectations, simultaneously confirms the deeper assumption of cosmic time in the standard cosmological model, contrary to the principles of GR. This empirical result directly motivates the CR framework, which preserves the empirical success of GR's space-time description while augmenting the mathematical structure in two key ways. First, the observed requirement for an absolute cosmic time motivates an explicit cosmological foliation of space-time. Second, the observational evidence for gravitational lensing and wave phenomena, which demonstrate that space dynamically warps, necessitates a diffeomorphic flexibility between foliated spatial slices and the real, evolving universe. These two principles distinguish CR from GR, as both a geometric augmentation of GR's ADM formalism and physically distinct theory: space-time is no longer treated as fundamental reality but as a phenomenological description of events unfolding in a physically real three-dimensional universe over cosmic time.

This refinement follows the natural progression of scientific theory, in which naive and superficial models that conflate phenomena with noumena give way to deeper, more coherent descriptions once empirical evidence exposes their limitations. The framework of CR emerges from empirical evidence, and naturally resolves the logical inconsistencies and predictive failures that arise when OR is extended to cosmology. The four examples explored in \S~\ref{sec:app} illustrate the significance of this distinction and highlight the explanatory power of the CR framework.

The first example demonstrated that the additional structural flexibility introduced by CR---diffeomorphism invariance between real space and cosmological spatial slices---provides a more natural setting for relativistic wave physics. Gravitational waves, electromagnetic waves, and other 3+1-dimensional wave phenomena fit coherently within this extended framework, which accommodates the physical evolution of waves within a dynamically evolving space rather than a rigid four-dimensional space-time manifold. This resolves an artificial constraint imposed by OR, where waves are treated as abstract perturbations of space-time rather than physical entities propagating through an evolving universe.

The second example considered the concept of time travel. CR explicitly prohibits closed timelike curves by construction, since the past is never physically real---only the present three-dimensional universe exists at any given moment of cosmic time. This simple but fundamental shift in ontology immediately resolves the paradoxes of time travel, making them not merely unlikely or prohibited, but categorically impossible. Unlike standard GR, which mathematically admits solutions with closed timelike curves, CR inherently forbids such structures as unphysical, ensuring logical consistency with our observed universe.

The third example addressed gravitational collapse, showing that CR leads to a fundamentally different description of black holes. While standard GR treats black holes as already formed objects with event horizons and singularities that have formed in finite time, CR demonstrates that collapse is asymptotic in cosmic time. Infalling matter never reaches the Schwarzschild radius within finite cosmic time, meaning that black holes in our universe are better understood as regions of continuous, asymptotic gravitational collapse rather than completed objects with well-defined interiors. This interpretation naturally resolves black hole paradoxes such as information loss, as no completed event horizon ever forms. This leads to a key empirical test: gravitational wave signals from black hole mergers should exhibit subtle deviations from the standard models, reflecting the fact that merging black holes are dynamically evolving regions of asymptotic collapse rather than pre-existing singularities.

Finally, the fourth example extended the framework beyond Euclidean cosmological CR descriptions, to construct a cosmological space-time model consistent with observed large-scale expansion. The SdS cosmology derived in this work followed a redefinition of causal structure within foliated dS space-time, preserving the absolute cosmic foliation. This methodological approach, which CR naturally accommodates within its basic differential geometric framework, leads to resolutions of all six epistemic problems of FLRW cosmology: (1) The cosmic time and isotropy problems are identically resolved through the CR framework itself. (2) The expansion problem is naturally incorporated as a feature of the basic cosmological dS framework, which describes a universe that must expand rather than one whose expansion contravenes the natural state of a universe filled with matter and radiation. (3) The horizon problem is resolved as uniform expansion is a geometric property of space-time rather than an emergent effect of matter distribution. (4) The flatness problem is eliminated as the phenomenological expansion rate of this model is identical to the flat $\Lambda$CDM expansion rate. (5) The energy budget problem is naturally avoided since expansion does not depend on the relative contributions of matter and radiation, including the predicted vacuum energy. 

The SdS model also provides an explanation for the observed large-scale uniformity of expansion. The predicted redshift contrast of $\sim1\%$ in the CMB due to expansion anisotropy is a large-scale statistical average and does not imply uniformity on local scales. In fact, general relativistic expectations suggest significant deviations in local expansion rates in overdense and underdense regions \cite{Wiltshire2009,Seifert2025}. However, observational evidence from cosmic flow studies \cite{Tully2014,Valade2024}, which subtract uniform expansion to reveal residual structure, presents a different picture: rather than large deviations, the residual velocity fields show coherent basins of attraction, consistent with the presence of an expansion threshold. This reinforces the conclusion that expansion is not merely uniform on average, but is an exact constraint imposed by the fundamental geometry, supporting the hypothesis of a threshold density above which structures cease to expand.

The SdS model also makes a critical prediction: the lack of an early radiation-dominated era influencing the expansion rate. This presents an opportunity to resolve the Hubble tension, as the standard approach relies on an FLRW-based mapping of early-universe physics to late-time observations. In this framework, careful treatment is required to map the evolution of the photon-baryon fluid from the noumenal photon frame to phenomenological appearances in the matter frame. This distinction must be accounted for in future work, particularly when determining the evolution of the sound horizon.

Furthermore, the SdS model reveals a tantalising possibility: the coordinate singularity at $r=0$ does not correspond to a physical singularity but instead marks a minimum radius, suggesting a cosmic bounce. This same structure describes a white hole expanding from $r=0$ at the beginning of cosmic time, naturally raising the question: could gravitational collapse in one universe seed the birth of another? In CR, black holes asymptotically collapse toward their event horizons in finite proper time but infinite cosmic time. This opens the possibility that the accumulated mass-energy of collapsing structures, reaching the event horizon at the end of cosmic time, could serve as the progenitor of a new universe. Remarkably, this interpretation is already empirically constrained: the phenomenological expansion history of the SdS model \textit{exactly} matches the observed expansion history of our universe. Had the observed expansion rate deviated from this form, this interpretation would be invalid. 

A direct observational test of CR lies in gravitational wave signals from binary black hole mergers. Since CR predicts that black holes are not pre-existing objects with fully formed event horizons but are instead dynamically evolving regions of asymptotic collapse, their ringdown waveforms should deviate from standard GR expectations. Specifically, post-merger signals should reflect the fact that event horizons never fully form in finite cosmic time, introducing subtle but measurable effects in the late-stage gravitational wave emission. Future high-precision gravitational wave observatories provide a means to directly test this distinction.

More speculatively, the SdS model suggests a novel possibility: the coordinate singularity at $r=0$ may not correspond to a physical singularity, but rather to a minimal radius, implying a cosmic bounce. This raises an intriguing hypothesis---could gravitational collapse in one universe seed the birth of another? If so, black holes may serve as bridges between cosmological cycles, with the accumulated mass-energy of collapsing structures reaching their event horizons at the end of cosmic time, triggering a transition into a new expansion phase. While this idea remains theoretical, its consequences are empirically testable: gravitational wave signatures from black hole mergers may encode residual imprints of a prior cosmic cycle. If deviations beyond those predicted by asymptotic collapse alone are observed, they could serve as indirect evidence of cosmic rebirth. 

If such additional deviations are detected, they would provide a direct observational probe of cosmic rebirth. Future high-precision gravitational wave observatories offer the opportunity to test this aspect of CR. More broadly, since CR predicts that black holes are dynamically evolving regions of asymptotic collapse rather than pre-existing singularities, gravitational wave observations of mergers serve as a crucial test. If black hole mergers exhibit deviations consistent with dynamically evolving horizons---and if infalling proper times remain finite at the horizon---this would indirectly support the hypothesis that gravitational collapse remains incomplete at the end of cosmic time, potentially leading to a cosmic bounce.

These empirical tests---including large-scale structure evolution and density thresholds, gravitational wave signatures, and a testable alternative resolution of the Hubble tension---demonstrate that CR is not merely a conceptual refinement of GR but a geometric augmentation directly necessitated by observational contradictions within OR, which leads to predictable outcomes in spite of its equivalent formulation of space-time. CR is not an arbitrary modification of GR---it is the necessary resolution to empirical failures that have persisted unrecognized within OR for a century. Its predictions are not speculative---they are the inevitable consequences of correcting an unexamined flaw in relativistic cosmology. Just as planetary epicycles were abandoned when the heliocentric model revealed a deeper geometric structure, the operationalist assumptions embedded in OR must now give way to a framework that directly aligns with empirical reality. The assumptions defining CR are empirically motivated, and its predictions offer clear paths to verification or falsification. The framework outlined in this work provides a coherent, logically consistent, and observationally grounded resolution to the fundamental problems in modern cosmology while offering unique solutions to paradoxes in time travel and black hole physics.

\bibliography{main}

\end{document}